% Generated draft: requires mathematical and source review.
\documentclass[11pt]{article}
\usepackage[margin=25mm]{geometry}
\usepackage{fontspec}
\setmainfont{Latin Modern Roman}
\usepackage{amsmath,amssymb,mathtools,unicode-math}
\setmathfont{Latin Modern Math}
\usepackage{enumitem,needspace,xurl,hyperref}
\setlength{\parindent}{0pt}
\setlength{\parskip}{5pt}
\setlength{\emergencystretch}{3em}
\newcommand{\N}{\mathbb{N}}
\newcommand{\Z}{\mathbb{Z}}
\newcommand{\Q}{\mathbb{Q}}
\newcommand{\R}{\mathbb{R}}
\newcommand{\C}{\mathbb{C}}
\newcommand{\F}{\mathbb{F}}
\newcommand{\A}{\mathbb{A}}
\newcommand{\PP}{\mathbb{P}}
\newcommand{\E}{\mathbb{E}}
\newcommand{\eps}{\varepsilon}
\newcommand{\dd}{\,\mathrm{d}}
\begin{document}
\section*{UnsolvedMath ID: 30000353}
\textbf{Source problem code:} catalogue entry
\subsection{Definitions and mathematical statement}
The record asks two questions about quasiperiodic Schr\"odinger cocycles. To explain a standard discrete model, let $V:\mathbb T^d\to\R$, let $\omega\in\R^d$ satisfy rational independence of $1,\omega_1,\ldots,\omega_d$, and write
\[
 (H_{\lambda,V,\theta}u)_j=u_{j+1}+u_{j-1}
       +\lambda V(\theta+j\omega)u_j,\qquad j\in\Z.
\]
The associated transfer matrix is
\[
 A_{E,\lambda}(\theta)=
 \begin{pmatrix}E-\lambda V(\theta)&-1\\1&0\end{pmatrix}.
\]
With $A^{(n)}(\theta)=A(\theta+(n-1)\omega)\cdots A(\theta)$, its Lyapunov exponent is, under the usual assumptions ensuring the limit,
\[
 L_\lambda(E)=\lim_{n\to\infty}\frac1n
     \int_{\mathbb T^d}\log\|A_{E,\lambda}^{(n)}(\theta)\|\,\dd\theta.
\]
This discrete formulation is explanatory: the source does not explicitly specify discrete versus continuous Schr\"odinger dynamics, base dimension, or its cocycle convention.

\begin{enumerate}
\item Let $E_0$ be the lowest spectral energy in the intended model, and assume $L_\lambda(E_0)>0$. Under sufficient regularity of the potential and rational independence of the frequency, determine the structure of the cocycle's minimal set; in particular, can its fibers be Cantor sets? In the usual projective-cocycle interpretation, the map acts on $\mathbb T^d\times\R\PP^1$ by
\[
 (\theta,[v])\longmapsto(\theta+\omega,[A_{E_0,\lambda}(\theta)v]).
\]
A minimal set is a nonempty compact invariant set containing no smaller nonempty compact invariant set; its fiber at $\theta$ is $\{[v]:(\theta,[v])\in K\}$. A Cantor fiber is a nonempty compact perfect totally disconnected fiber. The source does not identify a unique minimal set, its exact phase space, or whether Cantor fibers are asked for at one, every, or generic base point.
\item When the potential is only finitely smooth, does sufficiently large coupling force positive Lyapunov exponent for every energy? Once the intended potential and frequency class is fixed, the natural quantifier pattern is
\[
 \exists\lambda_0>0\ \forall\lambda\ge\lambda_0\ \forall E\in\R:
       L_\lambda(E)>0.
\]
The source does not specify whether large negative coupling is included, the differentiability order, potential nondegeneracy assumptions, stronger arithmetic frequency conditions, or the dependence of $\lambda_0$. No positive lower bound uniform in $E$ is added to pointwise positivity for all energies.
\end{enumerate}

The expressions ``sufficient regularity'' and ``finitely smooth'' do not determine admissible potential classes; in particular, they do not by themselves specify exclusions of constant or degenerate potentials. The questions must be read with missing report hypotheses still unresolved, rather than as a new universal theorem for every finitely smooth function. No analytic regularity or Diophantine hypothesis is silently imposed.
\subsection{Short English statement}
At the spectral bottom with positive Lyapunov exponent, the first question asks what minimal invariant sets look like and whether Cantor fibers occur. The second asks whether large coupling gives positivity at all energies for finitely smooth potentials. The precise model, regularity, and nondegeneracy hypotheses are missing.
\subsection{Sources}
\begin{itemize}
\item Local UnsolvedMath registry, permanent ID 30000353, code catalogue entry, ``Minimal Sets and Lyapunov Positivity for Quasiperiodic Schr\"odinger Cocycles.'' Proposed year 2005; proposer unspecified.
\item \emph{Dynamics of Cocycles and One-Dimensional Spectral Theory} (2006), Oberwolfach Reports, pp.\ 2933--2978, DOI 10.4171/owr/2005/51; \url{https://doi.org/10.4171/owr/2005/51}. The displayed report year 2006 differs from the DOI and proposed year 2005.
\item Supplied reference labeled ``Avila--Krikorian cocycle theorem'': \url{https://doi.org/10.4007/annals.2006.164.911}.
\item Statement status corrected\_verified; the supplied verification describes meaning-preserving editorial and LaTeX cleanup with no added source-dependent content.
\item Literature triage checked 2026-08-21: partially solved. It reports substantial analytic and special-subshift positivity and uniformity results, but says the general Cantor-fiber and finitely-smooth large-coupling questions remain open. Exact hypotheses and current status require independent verification. Created 2026-08-25; updated 2026-09-16.
\item Registry curation by Ulam AI / UnsolvedMath Contributors, CC BY 4.0; underlying report and literature provenance remain as supplied. Unreviewed statement-only draft.
\end{itemize}
\end{document}
